\section{Conclusion}
\label{sec:conclusion}

Most small decisions do not need a generative model. A typed question over
evidence, answered by an encoder with a law per type, gives a distribution a
caller can threshold, record and replay. Encoding options alone removes the
limits a shared sequence puts on option count, order and question count;
decoding a program's questions as typed variables, in waves bounded by depth
and projected onto rules, makes their answers coherent; joining every model
verdict with a deterministic policy makes safety the policy's, with the cost of
trusting the model counted exactly in shadow.

Training Kai taught as much as designing it. Local SGD across three kinds of
accelerator works and is exact to the bit, but an outer step that multiplies a
round's delta seven times, over a mixture this broad, cost the encoder its
ability to relate two texts while every single-text task improved; only
measurements aimed at the encoder itself found it. A barrier against the test
set is only as good as its notion of ``the same'': half of one baseline's test
set had same-wording siblings in its training data. Both are written down here
with the runs that show them.

Kai's case does not rest on winning every benchmark: Laya leads where it was
trained and trails Jev off its mixture, and Kai enters this paper only through
the tables its benchmark writes. It rests on what a hosted model cannot offer:
a pinned revision, deployment where the data live, no per-call fee, a runtime
linked into the system it governs, a package per run that says what was
decided, from what, and what would change it, and, in Enso, less generation per
successful task, which Section~\ref{sec:enso} is designed to measure.
